\documentclass[10pt,a4paper]{amsart}
\usepackage[margin=19mm]{geometry}
\usepackage{amsmath,amssymb,mathtools}
\usepackage[scr=rsfs,cal=euler]{mathalfa}
\usepackage{xfrac}
\usepackage[unicode,hidelinks,bookmarks=false]{hyperref}
\newcommand{\ie}{i.e.\ }
\newcommand{\cf}{cf.\ }
\newcommand{\resp}{respectively\ }
\newcommand{\Subsection}{\subsection}
\providecommand{\nobreakdash}{\nobreak\hbox{-}}
\setlength{\emergencystretch}{2em}
\multlinegap0pt
\usepackage{xcolor}
\usepackage{float}
\usepackage{amssymb}
\usepackage{algorithm}\usepackage{algpseudocode}
\usepackage[normalem]{ulem}
\usepackage{cancel}
\newtheorem{assumption}{Assumption}
\newtheorem{lemma}{Lemma}
\newtheorem{definition}{Definition}
\def\bbE{\mathbb E}
\def\bbR{\mathbb R}
\def\ind{\mathbb{I}}
\title{Generalized Regret Analysis of Thompson Sampling using Fractional Posteriors}
\author{Prateek Jaiswal, Debdeep Pati, Anirban Bhattacharya, Bani K. Mallick}
\date{Source: arXiv:2309.06349v2 (CC BY 4.0)}
\begin{document}
\maketitle

\noindent\emph{Source and licence.} Generalized Regret Analysis of Thompson Sampling using Fractional Posteriors. Version arXiv:2309.06349v2 (CC BY 4.0), distributed by arXiv under the Creative Commons Attribution 4.0 International licence, http://creativecommons.org/licenses/by/4.0/, as recorded in the arXiv licence metadata for this exact version and shown on the abstract page https://arxiv.org/abs/2309.06349v2. Full licence notice: \texttt{licenses/arxiv-2309.06349-CC-BY-4.0.txt}.\par\smallskip

\noindent\textit{Editorial scope.} Counted unit: the article's lemma lem:Post (source0.tex lines 590--605). The article's complete original proof of it is delivered with it (source0.tex lines 843--925, one \texttt{proof} environment of 83 lines, which the article places away from the statement). No mathematics is written, repaired or re-derived: the statement and the proof are the article's own lines, in the article's own notation, and no retained line is substituted or re-flowed.  Retained uncounted as the source-local context the proof needs: lines 378--381; lines 389--394; lines 678--681; lines 685--697; lines 734--739; lines 785--790; lines 819--826.  Nothing was omitted from any retained span, so the package carries no omission record.  No retained line contains a citation, so the wrapper carries no bibliography and no bibliography entry.  No retained line contains a figure environment, an image-inclusion command or drawing code, so no image is delivered and the package declares no asset dependency.  Editorial formatting only, in addition: an AMS \texttt{amsart} preamble that restates the article's own declarations and macros used by the retained lines, namely the environments \texttt{assumption}, \texttt{lemma}, \texttt{definition} on separate counters, together with the standard packages those lines need.  The retained environments print in this document as Assumption 1, Lemma 1, Definition 1 in order of appearance, whereas the article prints the same items with the numbering of its own full document.  The complete native source of the article as distributed in the arXiv e-print for this exact version is bundled unchanged as \texttt{sources/146921/source0.tex}.
      \begin{assumption}[Renyi divergence]~\label{ass:Renyi}
      For any $\alpha\in(0,1)$ and $\theta,\vartheta \in \Theta$, there exists a constant $D>0$, such that
     \(   |\bbE_{\theta}[r]- \bbE_{\vartheta}[r]| \leq D\sqrt{\frac{2}{\alpha} D_{\alpha}(\theta,\vartheta)},\) where $\bbE_{\theta}[r]$ is the mean reward under reward distribution $p(\cdot|\theta)$. 
\end{assumption}

\begin{assumption}[Prior thickness]~\label{ass:prior2}
            We assume that, for each $k\in[K]$, $\alpha\in(0,1)$, $j\geq 1$, and $\delta_k^2=\Delta_k^2/9$, for  a ball $B\left(\theta_{0,k}, C(\alpha)\delta_k^2 \right) := \left\{  \theta_k \in \Theta:D_2(\theta_{0,k},\theta_k)\leq  C(\alpha)\delta_k^2 \right\} $ and shell $F_k^j:= \big\{\theta_k \in \Theta: D\frac{\alpha}{2} (j+1) \delta_k^2 \geq D_{\alpha}(\theta_k,\theta_{0,k}) \geq D j \frac{\alpha}{2} \delta_k^2  \big\}$, there exist $\kappa > 0$, such that 
        \begin{align}\tag{B1} \frac{\Pi_k(F_k^j)}{\Pi_k(B\left(\theta_{0,k},C(\alpha)\delta_k^2 \right) )} \leq \kappa(j+1),
            \end{align}
        where $\kappa$ (a constant) depends on the prior, reward distributions and true model parameter.
        \end{assumption}

Below we specify two families of distribution that satisfies all the regularity conditions that we have used to derive the regret bounds.  The first family is the popular sub-Gaussian family of distributions. one specifies the regularity conditions for a sub-Gaussian reward. \begin{definition}[Sub-Gaussian]~\label{def:SubReward}
   For any arbitrary $\Theta$ and for each arm $k$, we assume that
    the reward distribution $P_{\theta_k}$ is sub-Gaussian with parameter $\sigma=1$, that is $\bbE_{P_{\theta_k}} [e^{s x_k}] \leq e^{s \theta_k + s^2/2}$ for all $s\in \bbR$.
    \end{definition}

    \begin{definition}[1-d Exponential]~\label{def:rewardExp}
        We assume for $\Theta \subseteq \bbR$ that 
        \begin{itemize}
            \item[i.] $p(x|\theta_k) = e^{x \theta_k   -  A(\theta) + C(x) }$, with mean as $ A'(\theta)$ and variance as $A''( \theta)$, where $A(\cdot)$ is the log-partition function and $C(\cdot)$ is some known mapping. We also define the link function $g(\cdot)=A'(\cdot)$.  
            \item[ii.] the link function is Lipschitz continuous, that is for any $x,y\in \Theta$, there exists a constant $C_g>0$, such that
             $   |g(x)-g(y)|\leq C_g|x-y|$, and
            \item[iii.] the log-partition function is strongly convex with parameter $m$, that is for any $\gamma\in (0,1)$, $x,y\in \Theta$, there exists an $m>0$, such that
            \begin{align*}
                \gamma A(x) + (1-\gamma)A(y ) - A(\gamma x + (1-\gamma)y ) 
                \geq \frac{1}{2} \gamma (1-\gamma) m |x-y|^2.
            \end{align*}
        \end{itemize}
    \end{definition} 

\begin{lemma}[R\'enyi transportation cost inequality]~\label{lem:Pin}
    Fix $\alpha\in (0,1)$. For any two sub-Gaussian measures $\mu$ and $\nu$ with sub-Gaussian parameter $\sigma_{\mu}$ and $\sigma_{\nu}$ respectively, such that $\nu$ is absolutely continuous wrt $\mu$, that is $\nu << \mu$, the $\alpha-$R\'enyi divergence can be bounded below by the absolute difference between the respective means. In particular, for any random variable $X$ having measure $\mu$ and $\nu$, we have  
    \begin{align}
        |\bbE_{\nu}[X]-\bbE_{\mu}[X]  | &\leq \sqrt{(\sigma_{\mu}^2\alpha+\sigma_{\nu}^2(1-\alpha))} \sqrt{\frac{2}{\alpha} D_{\alpha}(\nu\|\mu) }.
        \end{align} 
\end{lemma}

\begin{lemma}~\label{lem:Exp}
    Fix $\alpha\in (0,1)$. Let $r$ be a random variable having its distributions lying an exponentially family with parameters $\theta$ and $\theta_0$ but same $A(\cdot)$ and $C(\cdot)$. Then under Definition~\ref{def:rewardExp} the $\alpha-$R\'enyi divergence can be bounded below by the absolute difference between the respective means. In particular, we have  
    \begin{align}
        |\bbE_{\theta}[r]-\bbE_{\theta_0}[r]  | &\leq \frac{C_g}{\sqrt{m}} \sqrt{\frac{2}{\alpha} D_{\alpha}(\theta,\theta_0) }.
        \end{align} 
\end{lemma}

\begin{lemma}[Adaptively stopped likelihood ratios]\label{lem:MG}
Fix $k \in [K]$ and a \emph{fixed, nonrandom} parameter $\theta_k \in \Theta$. Let $g \geq 0$ be a measurable function on the reward space with $m_g(\theta_k) := \bbE_{x \sim P_{\theta_{0,k}}}[g(x)] \in (0,\infty)$; note that $m_g(\theta_k)$ is a deterministic constant. Then, for any bounded stopping time $\tau \leq T$ with respect to $(\mathcal{F}_s)$, any deterministic $\Lambda \geq 0$, and any $\rho > 0$ with $\rho\, m_g(\theta_k) \leq 1$,
\begin{align}
    \bbE\left[ \rho^{\bar n_k(\tau)} \prod_{i=1}^{\bar n_k(\tau)} g(x_{i,k})\, \ind\{\bar n_k(\tau) \geq \Lambda\} \right] \leq \big(\rho\, m_g(\theta_k)\big)^{\Lambda}.
    \label{eq:MGkey}
\end{align}

\end{lemma}

        \begin{lemma}[$\alpha$-posterior concentration]~\label{lem:Post}
            Fix {$k\in [K]$,} and $\alpha\in(0,1)$ . Under  Assumptions~\ref{ass:Renyi} and~\ref{ass:prior2}, for  any bounded stopping time $\tau \leq T$ with respect to the filtration $\mathcal{F}_s := \sigma\left(X_s, \{i_r\}_{r=1}^s\right)$, provided $C(\alpha)\delta_k^2 n_{k,0} \geq 1$ (this condition is needed only for the first display below)  and $\kappa \leq (1-e^{-4})^2/48$, 
            \begin{align}
                \bbE& \left[ \bbE\left[ \ind\{|\mu_k^{\tau} -\mu_{0,k} |^2 \geq \delta_k^2 \} | X_{\tau} \right] e^{C(\alpha) \bar n_k(\tau) \delta_k^2} \right] 
                \leq 2^{-1},
                \label{eq:postcon}
            \end{align}
            Moreover, for any deterministic $\Lambda$ satisfying $C(\alpha)\delta_k^2 \Lambda \geq 1$, we also have the indicator form, which does \emph{not} require the condition $C(\alpha)\delta_k^2 n_{k,0} \geq 1$:
            \begin{align}
                \bbE \left[ \bbE\left[ \ind\{|\mu_k^{\tau} -\mu_{0,k} |^2 \geq \delta_k^2 \} | X_{\tau} \right] \ind\{\bar n_k(\tau) \geq \Lambda\} \right] 
                \leq 2^{-1} e^{ - C(\alpha) \Lambda \delta_k^2 },
                \label{eq:postconL}
            \end{align}
            where $\bar n_k(\tau) := n_k(\tau)+n_{k,0}$ counts the pulls of arm $k$ up to time $\tau$ including the warm-up samples, $\mu_k^{\tau}$ denotes the mean corresponding to one draw from the $\alpha$-posterior $\pi_{\alpha}(\cdot|X_{\tau})$ %(so that the inner expectation equals the $\alpha$-posterior mass of the displayed event; for a deterministic $\tau = t$ this is exactly the quantity $\mu_k^t$ used elsewhere)
            , and $C(\alpha) = D (1-\alpha) \min \left(2\alpha,1-\alpha  \right)/16$ .
        \end{lemma}

\begin{proof}[Proof of Lemma~\ref{lem:Post}]
    Using Lemma~\ref{lem:Pin} for sub-Gaussian rewards satisfying conditions in Definition~\ref{def:SubReward}   and Lemma~\ref{lem:Exp} for Exponential family rewards satisfying Definition~\ref{def:rewardExp}, we have
    \begin{align}
        \left|   \mu_k-\mu_{0,k}  \right| %=\left|\int_{[0,1]} [r p(r|\theta_{0,k}) - r p(r|\theta_k) ]dr  \right| 
        \leq \frac{1}{\sqrt{D}}\sqrt{\frac{2D_{\alpha}(\theta_k,\theta_{0,k}) }{\alpha}},
        \label{eq:Pins}
        \end{align}
    where $D=1$ for sub-Gaussian rewards and $D=\frac{m}{C_g^2}$ for exponential family rewards. We omit using superscript $t$ in $\mu_k$ and $\theta_k$ for brevity.
    Therefore, it follows from the inequality above that 
    \begin{align}
        E\left[ \ind\{|\mu_k -\mu_{0,k} |^2 \geq \delta_k^2 \} | X_{t} \right] \leq E\left[ \ind\{D_{\alpha}(\theta_k,\theta_{0,k}) \geq D \frac{\alpha}{2} \delta_k^2  \} | X_{t} \right].
        \label{eq:MABDE2}
    \end{align} 
   Now for any $k\in[K]$, let us define a set 
    \begin{align}
        F_k = \left\{\theta_k \in \Theta: D_{\alpha}(\theta_k,\theta_{0,k}) \geq D \frac{\alpha}{2} \delta_k^2  \right\}.
        \label{eq:MABDE2b}
    \end{align}

 and define $\bar n_k(t)=n_k(t)+n_{k,0} =\sum_{s=1}^{t}  \ind(i_s = k) + n_{k,0} $. Observe that 
\begin{align}
    \nonumber
    \Pi_{\alpha}(F_k|X_{t})&= \frac{\int_{F_k} \pi_k(\theta_k) p(X_{k,0}|\theta_{k})^{\alpha}\prod_{s=1}^{t} p(x_{i_s}|\theta_{i_s})^{\alpha \ind(i_s = k)} d\theta_k} {\int \pi_k(\theta_k) p(X_{k,0}|\theta_{k})^{\alpha}\prod_{s=1}^{t} p(x_{i_s} |\theta_{i_s})^{\alpha \ind(i_s = k)} d\theta_k }
        \\
        \nonumber
        &= \frac{\int_{F_k} \pi_k(\theta_k) \prod_{i=1}^{\bar n_k(t)} p(x_{i,k}|\theta_k)^{\alpha }d\theta_k} {\int \pi_k(\theta_k) \prod_{i=1}^{\bar n_k(t)} p(x_{i,k}|\theta_k)^{\alpha }d\theta_k }
        \\
        &= \frac{\int_{F_k} \pi_k(\theta_k) e^{-\alpha \ell_{\bar n_k(t)}(\theta_k,\theta_{0,k})}d\theta_k} {\int \pi_k(\theta_k) e^{-\alpha \ell_{\bar n_k(t)}(\theta_k,\theta_{0,k})}d\theta_k },
        \label{eq:MABDE3}
\end{align}
where $\ell_{\bar n_k(t)}(\theta_k,\theta_{0,k})= \log \frac{\prod_{i=1}^{\bar n_k(t)} p(x_{i,k}|\theta_{0,k})}{\prod_{i=1}^{\bar n_k(t)} p(x_{i,k}|\theta_k)}$.
Next define a set 
\[A_{\bar n_k(t)} = \left\{ X_{k,t}: \int_{\Theta} \tilde \pi_k(\theta_k) e^{-\alpha \ell_{\bar n_k(t)}(\theta_k,\theta_{0,k})}d\theta_k \leq 4^{-\alpha} e^{-D \frac{(\alpha-\alpha^2)}{8} \bar n_k(t)\delta_k^2}  \right\},\] where for any $B\subset \Theta$, $\tilde \Pi_k(B) = \Pi_k\left(B\cap \Theta\right)/\Pi_k\left(B\left(\theta_{0,k},\delta_k \right)\right)$.
Fix a bounded stopping time $\tau \leq T$ with respect to the filtration $\mathcal{F}_s := \sigma\left(X_s, \{i_r\}_{r=1}^s\right)$; recall that $C(\alpha)\delta_k^2 n_{k,0} \geq 1$ by assumption. We abbreviate $\bar n_k := \bar n_k(\tau)$; all pathwise identities above are evaluated at time $\tau$. Our only probabilistic tool is the optional-stopping bound of Lemma~\ref{lem:MG}. %We stress that in every application below $\theta_k$ is a fixed, nonrandom parameter: it plays the role of the variable of integration under Fubini--Tonelli, never of the posterior draw. The posterior draw $\mu_k^{\tau}$ enters only through the pathwise bound~\eqref{eq:MABDE2}.
In view of~\eqref{eq:MABDE2} and~\eqref{eq:MABDE3}, it suffices to bound the compensated posterior mass
\begin{align}
     \bbE\left[\Pi_{\alpha}(F_k|X_{\tau})\, e^{C(\alpha) \bar n_k \delta_k^2}\right] &\leq \bbE\left[\ind\big(A_{\bar n_k}\big)\, e^{C(\alpha) \bar n_k \delta_k^2}\right] \nonumber
     \\
     &\quad + \bbE \left[  \frac{\int_{F_k} \pi_k(\theta_k) e^{-\alpha \ell_{\bar n_k}(\theta_k,\theta_{0,k})}d\theta_k} {\int \pi_k(\theta_k) e^{-\alpha \ell_{\bar n_k}(\theta_k,\theta_{0,k})}d\theta_k } \ind\big(A_{\bar n_k}^\mathtt{c}\big)\, e^{C(\alpha) \bar n_k \delta_k^2} \right].
    \label{eq:MABDE4}
\end{align}

\emph{Second term.} On the set $A_{\bar n_k}^\mathtt{c}$, the definitions of $A_{\bar n_k}$ and of $\tilde \pi_k$ give the pathwise lower bound $\int \pi_k e^{-\alpha \ell_{\bar n_k}} \geq \Pi_k(B\left(\theta_{0,k},\delta_k \right))\, 4^{-\alpha} e^{-D \frac{(\alpha-\alpha^2)}{8} \bar n_k \delta_k^2}$ on the denominator. Writing $F_k= \bigcup_{j\geq 1} F_k^j$ with $F_k^j := \left\{\theta_k \in \Theta: D\frac{\alpha}{2} (j+1) \delta_k^2 \geq D_{\alpha}(\theta_k,\theta_{0,k}) \geq D j \frac{\alpha}{2} \delta_k^2  \right\}$, we obtain the pathwise bound
\begin{align}
    &\frac{\int_{F_k} \pi_k e^{-\alpha \ell_{\bar n_k}}} {\int \pi_k e^{-\alpha \ell_{\bar n_k}} } \ind\big(A_{\bar n_k}^\mathtt{c}\big)\, e^{C(\alpha) \bar n_k \delta_k^2} \nonumber \\
    &\qquad \leq \frac{4^{\alpha}}{\Pi_k(B\left(\theta_{0,k},\delta_k \right))} \sum_{j \geq 1} \int_{F_k^j} \pi_k(\theta_k)\, e^{-\alpha \ell_{\bar n_k}(\theta_k,\theta_{0,k})}\, e^{\left(D \frac{(\alpha-\alpha^2)}{8} + C(\alpha)\right) \bar n_k \delta_k^2}\, d\theta_k.
    \label{eq:MABDE6}
\end{align}
For a fixed $\theta_k$, note that $e^{-\alpha \ell_{\bar n_k}(\theta_k,\theta_{0,k})} = \prod_{i=1}^{\bar n_k} g_1(x_{i,k})$ for $g_1(x) = \left( p(x|\theta_k)/p(x|\theta_{0,k})\right)^{\alpha}$, and
\begin{align}
    m_{g_1}(\theta_k) = \int p(x|\theta_k)^{\alpha}\, p(x|\theta_{0,k})^{1-\alpha}\, dx = e^{-(1-\alpha) D_{\alpha}(\theta_k,\theta_{0,k})}.
    \label{eq:Iden1}
\end{align}
On $F_k^j$ we have $D_{\alpha}(\theta_k,\theta_{0,k}) \geq D j \frac{\alpha}{2}\delta_k^2$, so taking $\rho = e^{\left(D\frac{(\alpha-\alpha^2)}{8} + C(\alpha)\right)\delta_k^2}$ and using $C(\alpha) \leq D\frac{(\alpha-\alpha^2)}{8}$ gives $\rho\, m_{g_1}(\theta_k) \leq e^{-D \frac{(4j-2)(\alpha-\alpha^2)}{8}\delta_k^2} \leq 1$. Applying Lemma~\ref{lem:MG} (with $g = g_1$, this $\rho$, and $\Lambda = n_{k,0}$, whose indicator is vacuous) for each fixed $\theta_k \in F_k^j$ and Fubini--Tonelli,
\begin{align}
    \bbE\left[ \int_{F_k^j} \pi_k(\theta_k)\, e^{-\alpha \ell_{\bar n_k}(\theta_k,\theta_{0,k})}\, e^{\left(D \frac{(\alpha-\alpha^2)}{8} + C(\alpha)\right) \bar n_k \delta_k^2}\, d\theta_k\right] \leq \Pi_k(F_k^j)\, e^{-D \frac{(4j-2)(\alpha-\alpha^2)}{8}\, n_{k,0}\, \delta_k^2 }.
    \label{eq:MABDE6.5}
\end{align}
Therefore, using Assumption~\ref{ass:prior2} to bound the ratios $\Pi_k(F_k^j)/\Pi_k(B\left(\theta_{0,k},\delta_k \right)) \leq \kappa (j+1)$, the expectation of the second term in~\eqref{eq:MABDE4} is bounded by
\begin{align}
    \kappa\, 4^{\alpha}\sum_{j\geq 1} (j+1)\, \xi^{4j-2} \leq \frac{3\kappa 4^{\alpha}}{(1-e^{-4})^2}\, \xi^{2} \leq 4^{-1}.
    \label{eq:MABDE7}
\end{align}
Here $\xi := e^{-D \frac{(\alpha-\alpha^2)}{8}\, n_{k,0}\, \delta_k^2 }$, the series computation $\sum_{j\geq 1}(j+1)\xi^{4j-2} \leq 3 \xi^2/(1-\xi^4)^2$ is elementary, the bound $\xi \leq e^{-C(\alpha) n_{k,0} \delta_k^2} \leq e^{-1}$ holds because $C(\alpha) \leq D(\alpha-\alpha^2)/8$ and $C(\alpha) n_{k,0} \delta_k^2 \geq 1$, and the last inequality uses the assumption $3\kappa 4/(1-e^{-4})^2 \leq 4^{-1}$.

\emph{First term.} On the set $A_{\bar n_k}$ we have, pathwise, $\left[ \int \tilde \pi_k(\theta_k) e^{-\alpha \ell_{\bar n_k}(\theta_k,\theta_{0,k})} d\theta_k \right]^{-1/\alpha} \geq 4\, e^{ D \frac{1-\alpha}{8} \bar n_k \delta_k^2 }$, and hence, using Jensen's inequality for the convex map $x \mapsto x^{-1/\alpha}$,
\begin{align}
    \ind\big(A_{\bar n_k}\big)\, e^{C(\alpha) \bar n_k \delta_k^2} \leq 4^{-1} \int \tilde \pi_k(\theta_k)\, e^{\ell_{\bar n_k}(\theta_k,\theta_{0,k})}\, e^{ -\left(D \frac{1-\alpha}{8} - C(\alpha)\right) \bar n_k \delta_k^2 }\, d\theta_k.
    \label{eq:MABDE8}
\end{align}
For fixed $\theta_k$, $e^{\ell_{\bar n_k}(\theta_k,\theta_{0,k})} = \prod_{i=1}^{\bar n_k} g_2(x_{i,k})$ with $g_2(x) = p(x|\theta_{0,k})/p(x|\theta_k)$ and $m_{g_2}(\theta_k) = e^{D_2 \left( {\theta_{0,k}}\| {\theta_k} \right)}$. On the support of $\tilde \pi_k$ we have $D_2 \left( {\theta_{0,k}}\| {\theta_k} \right) \leq D \frac{(\alpha-\alpha^2)}{8} \delta_k^2$ by the definition of the set $B\left(\theta_{0,k},\delta_k \right)$, so with $\rho = e^{-\left(D \frac{1-\alpha}{8} - C(\alpha)\right)\delta_k^2}$ and $C(\alpha) \leq D\frac{(1-\alpha)^2}{16}$ we get $\rho\, m_{g_2}(\theta_k) \leq e^{-D \frac{(1-\alpha)^2}{16}\delta_k^2} \leq 1$. Applying Lemma~\ref{lem:MG} (with $g = g_2$ and $\Lambda = n_{k,0}$) and Fubini--Tonelli once more,
\begin{align}
    \bbE\left[\ind\big(A_{\bar n_k}\big)\, e^{C(\alpha) \bar n_k \delta_k^2}\right] \leq 4^{-1} \big(\rho\, m_{g_2}(\theta_k)\big)^{n_{k,0}} \leq 4^{-1}.
    \label{eq:MABDE8b}
\end{align}
Combining~\eqref{eq:MABDE7} and~\eqref{eq:MABDE8b} in~\eqref{eq:MABDE4},
\begin{align}
    \bbE\left[\Pi_{\alpha}(F_k|X_{\tau})\, e^{C(\alpha) \bar n_k \delta_k^2}\right]
    \leq 4^{-1} + 4^{-1} = 2^{-1},
    \label{eq:MABDE9}
\end{align}
which, together with~\eqref{eq:MABDE2}, proves~\eqref{eq:postcon}. The indicator form~\eqref{eq:postconL} holds under its own hypothesis $C(\alpha) \Lambda \delta_k^2 \geq 1$ alone, without requiring $C(\alpha) n_{k,0} \delta_k^2 \geq 1$: the identical argument applies with the compensation factor $e^{C(\alpha)\bar n_k \delta_k^2}$ replaced throughout by the indicator $\ind\{\bar n_k \geq \Lambda\}$ (so the choices of $\rho$ lose their $C(\alpha)$ terms) and with Lemma~\ref{lem:MG} invoked at threshold $\Lambda$ in place of $n_{k,0}$; every estimate then holds with $\Lambda$ in place of $n_{k,0}$, and the series bound uses $\xi = e^{-D\frac{(\alpha-\alpha^2)}{8} \Lambda \delta_k^2} \leq e^{-C(\alpha) \Lambda \delta_k^2} \leq e^{-1}$, yielding $\bbE[\Pi_{\alpha}(F_k|X_{\tau})\, \ind\{\bar n_k \geq \Lambda\}] \leq 2^{-1} e^{-2C(\alpha) \Lambda \delta_k^2} \leq 2^{-1} e^{-C(\alpha) \Lambda \delta_k^2}$.
\end{proof}

\end{document}
